\documentclass[11pt]{amsart}
\usepackage[T1]{fontenc}
\usepackage{lmodern}
\usepackage{amsmath,amssymb,amsthm}
\usepackage[margin=1.1in]{geometry}
\usepackage{booktabs}
\usepackage{hyperref}

\newtheorem{theorem}{Theorem}
\newtheorem{lemma}[theorem]{Lemma}
\newtheorem{proposition}[theorem]{Proposition}
\newtheorem{corollary}[theorem]{Corollary}
\newtheorem{conjecture}[theorem]{Conjecture}
\newtheorem{question}[theorem]{Question}
\theoremstyle{definition}
\newtheorem{definition}[theorem]{Definition}
\newtheorem{remark}[theorem]{Remark}
\newtheorem{example}[theorem]{Example}
\newtheorem*{remark*}{Remark}
\newtheorem*{example*}{Example}

\newcommand{\Ev}{\mathrm{Ev}}
\newcommand{\Od}{\mathrm{Od}}
\newcommand{\cF}{\mathcal F}
\newcommand{\cP}{\mathcal P}

\title[Prescribed ends in path decompositions]{Prescribed ends in path decompositions: a cut inequality, cut selection, and hardness}
\author{Justin Williams}
\date{November 2026 (working draft)}

\begin{document}
\begin{abstract}
Let $G$ be a connected graph of even order $n=2k$, and let $M$ consist of half of its vertices of even
degree. The \emph{simple profile} of $M$ asks for a partition of $E(G)$ into $k$ paths with one end at
every odd vertex, two ends at every vertex of $M$, and none at the other even vertices. Such a
partition attains Gallai's bound. A cut inequality of the even-degrees manuscript is necessary for
it. We study this inequality.
\begin{itemize}
\item \emph{Equivalent forms.} The inequality has four equivalent forms. The most useful asks that
the edges at each vertex $v$ extend to edge-disjoint \emph{arms} in $G-v$ ending at the prescribed
ends. For half-sets it becomes a single-commodity cut condition in the graph $G^v$ obtained by
complementing the star at $v$.
\item \emph{Cut selection.} The inequality can always be met. Every connected graph of even order has
an explicit pairing of its even vertices such that every half-set with one vertex from each pair
satisfies all the inequalities. The construction takes the $T$-join of a spanning tree. At each
vertex it chooses the transitions of the $T$-join by the $T$-join pairing lemma in the free part of
$G^v$. A parity lemma makes the odd vertices there exactly the ends of these transitions.
\item \emph{Sufficiency fails, and characterisation is hard.} The inequality suffices for every graph
of order at most eight and in the dense range, but not in general. Following a construction of the
project lead, we prove the following. For every connected bridgeless cubic graph $H$ there is a graph
$G(H)$, with degrees $2$ and $4$, and a half-set of it that satisfies every cut inequality and is
realisable if and only if $H$ is Hamiltonian. The proof combines Kotzig's theorem on Hamiltonian
decompositions of line graphs with a computation of the cuts of line graphs. Hence deciding
realisability of a half-set is NP-complete, even for half-sets that meet every cut inequality.
\item \emph{Matroid structure.} For each vertex, the half-sets meeting its inequalities are the bases
of a matroid. The cut-feasible family is the intersection of these matroids, and is not a matroid in
general. In cones it is described exactly by edge-disjoint pairings avoiding the apex.
\end{itemize}
We also restate several results of the archive as cut statements and discuss related work.
\end{abstract}
\maketitle

\section{Setting}\label{sec:setting}

All graphs are finite and simple. A path has at least one edge, and its ends are its two distinct end
vertices. A \emph{partition} of a graph is a partition of its edge set into paths. Throughout, $G$ is
connected of even order $n=2k\ge 2$. We write $d(v)$ for degree, $N(v)$ for the neighbourhood and
$e(v,Y)=|N(v)\cap Y|$. $\Od$ and $\Ev$ are the sets of odd and even vertices, and $e=|\Ev|$, which is
even. Since $G$ is connected, every even vertex has degree at least two. For $Y\subseteq V$,
$\delta(Y)$ is the edge cut. For $v\notin Y$, $\delta_{G-v}(Y)$ is the set of edges between $Y$ and
$V-Y-v$.

A \emph{profile} is a map $b:V\to\mathbb Z_{\ge 0}$, and $b(Y)=\sum_{y\in Y}b(y)$. A partition
\emph{realises} $b$ if it has exactly $b(v)$ path ends at each vertex $v$. A realisation has $b(V)/2$
paths. A \emph{half-set} is a set $M\subseteq\Ev$ with $|M|=e/2$. Its \emph{simple profile} is
\[
b_M=1 \text{ on }\Od,\qquad b_M=2\text{ on }M,\qquad b_M=0\text{ on }\Ev\setminus M,
\]
so $b_M(V)=n$ and a realisation of $b_M$ has exactly $k$ paths. Thus Gallai's conjecture for $G$
follows if some half-set is realisable. The \emph{local conditions} for a profile with $b(V)=2k$ are
\[
b(v)\equiv d(v)\pmod 2,\qquad b(v)\le d(v),\qquad d(v)+b(v)\le 2k .
\]
They are necessary, and they hold automatically for simple profiles: $d(v)\le n-1$, and
$d(v)\le n-2$ when $d(v)$ is even.

\emph{Overview.} Section~\ref{sec:emm} gives the cut inequality and its forms. Section~\ref{sec:consequences}
sets up the selection statements and the matroid structure. Section~\ref{sec:cs} proves cut selection
(Theorem~\ref{thm:ps}), and Section~\ref{sec:hard} shows that the inequality is not sufficient and
that realisability is NP-complete (Theorem~\ref{thm:hard}). Sections~\ref{sec:dictionary}--\ref{sec:evidence}
contain the dictionary, related work and computations.

We cite the archive by the abbreviations of \texttt{research/synthesis.md}:
ED (even-degrees), SP (separating-pairs), TVS (two-vertex-separators), RE (root-endpoint-costs),
MR (marked-reduction), ES (endpoint-structure) and JM-repair.

\section{The cut inequality}\label{sec:emm}

ED~\S12.4 states the following inequality for a profile $b$ with $b(V)=2k$, a vertex $v$ and a set
$W\subseteq V-v$:
\begin{equation}\label{eq:one}
b(W)\le |\delta(W)|+2k-d(v)-b(v).
\end{equation}

\begin{lemma}[four forms]\label{lem:forms}
Let $b$ be a profile with $b(V)=2k$ and let $v$ be a vertex. The following are equivalent.
\begin{enumerate}
\item[(a)] Inequality \eqref{eq:one} holds for every $W\subseteq V-v$.
\item[(b)] $d(v)+b(v)\le b(U)+|\delta(U)|$ for every $U\ni v$ (demand form).
\item[(c)] $e(v,Y)\le b(Y)+|\delta_{G-v}(Y)|$ for every $Y\subseteq V-v$ (arm form).
\item[(d)] There are paths $A_y$ $(y\in N(v))$ of $G-v$ with the following properties. They are
pairwise edge-disjoint. $A_y$ starts at $y$ and may have length $0$. Every vertex $z\ne v$ is the
last vertex of at most $b(z)$ of them. (An \emph{arm system} at $v$.)
\end{enumerate}
\end{lemma}

\begin{proof}
(a)$\Leftrightarrow$(b): put $U=V-W$. Then $b(W)=2k-b(U)$ and $\delta(W)=\delta(U)$, so
\eqref{eq:one} becomes $d(v)+b(v)\le b(U)+|\delta(U)|$. The sets $W\subseteq V-v$ correspond exactly
to the sets $U\ni v$.

(b)$\Leftrightarrow$(c): put $U=Y\cup\{v\}$. Then $b(U)=b(v)+b(Y)$ and
$|\delta(U)|=|\delta_{G-v}(Y)|+d(v)-e(v,Y)$. Substituting gives (c).

(c)$\Leftrightarrow$(d): build a network from $G-v$. Every edge becomes two opposite arcs of
capacity $1$. Add a source $\sigma$ with an arc of capacity $1$ to each $y\in N(v)$, and a sink $\tau$
with an arc of capacity $b(z)$ from each $z\ne v$. An arm system is the same thing as an integral flow
of value $d(v)$. Indeed, cancel opposite flows on each edge, discard cycles and decompose the rest
into $\sigma$--$\tau$ paths. A cut is given by the set $X\subseteq V-v$ of vertices on the source
side, and its capacity is $|N(v)\setminus X|+|\delta_{G-v}(X)|+b(X)$. By the max-flow min-cut
theorem, a flow of value $d(v)$ exists if and only if every such capacity is at least $d(v)$, and that
is (c).
\end{proof}

\begin{proposition}[necessity; ED \S12.4]\label{prop:nec}
If a partition $\cP$ realises $b$, then (a)--(d) hold at every vertex.
\end{proposition}

\begin{proof}
We construct an arm system at $v$. Each edge $vy$ lies on a unique path $Q\in\cP$. Traverse $Q$ from
$v$ through $y$ to the end of $Q$ on that side. Let $A_y$ be this walk with the edge $vy$ removed. It
is a path of $G-v$ from $y$ to an end $z\neq v$ of $Q$, since $Q$ visits $v$ only once. Two arms
from different paths are edge-disjoint. The two arms of a path through $v$ are edge-disjoint, and they
end at its two different ends. Hence every end occurrence of every path terminates at most one arm,
and $z$ is the last vertex of at most $b(z)$ arms.
\end{proof}

ED proves the necessity of~\eqref{eq:one} by the same count. Each path through $v$ gives each of its
ends in $W$ its own edge of $\delta(W)$, and each of the $k-(d(v)+b(v))/2$ paths missing $v$ has at
most two ends in $W$.

\begin{definition}
A half-set $M$ is \emph{cut-feasible} if $b_M$ satisfies (a)--(d) at every vertex.
\end{definition}

\begin{question}[EMM, endpoint min-max]\label{conj:emm}
Is every cut-feasible half-set realisable?
\end{question}

For $W=\emptyset$, inequality~\eqref{eq:one} is the local condition $d(v)+b(v)\le 2k$. For $e=0$ the
empty half-set is trivially cut-feasible, and its realisability is Lov\'asz's theorem for graphs with
all degrees odd~\cite{L68}. The answer to Question~\ref{conj:emm} is yes for every connected graph of
order at most $8$ (Table~\ref{tab:emm}) and in the dense range (Proposition~\ref{prop:dense} with SP
Theorem~14.1). It is no in general (Section~\ref{sec:hard}).

\begin{remark}[Eulerian form]\label{rem:euler}
Let $G_b$ be $G$ together with a new vertex $t$ joined to each $z$ by $b(z)$ parallel edges. For a
simple profile, $G_b$ is Eulerian. Form (b) at every $v$ says $\lambda_{G_b}(v,t)=\deg_{G_b}(v)$ for
every $v\in V$: every vertex is joined to $t$ by as many edge-disjoint paths as its degree. A
realisation is a decomposition of $G_b$ into $k$ cycles through $t$, none of length two.
Question~\ref{conj:emm} asks whether this connectivity condition suffices for such special Eulerian
graphs; Section~\ref{sec:hard} shows that it does not. The closest
classical min-max theorem of this shape is the Lov\'asz--Cherkassky theorem on edge-disjoint
$T$-paths in inner-Eulerian graphs (\texttt{research/synthesis.md}~\S7).
\end{remark}

\subsection{The signed form for half-sets}
For a half-set $M$, define
\[
x_M=+1 \text{ on } \Ev\setminus M,\qquad x_M=-1\text{ on }M,\qquad x_M=0\text{ on }\Od.
\]
For a vertex $v$, let $G^v$ be the graph obtained from $G$ by deleting the edges at $v$ and joining
$v$ to every vertex of $V-v-N(v)$. For $Y\subseteq V-v$,
\[
|\delta_{G^v}(Y)|=|\delta_{G-v}(Y)|+|Y\setminus N(v)|.
\]

\begin{lemma}[signed form]\label{lem:signed}
A half-set $M$ satisfies (c) at $v$ if and only if
\begin{equation}\label{eq:signed}
x_M(Y)\le|\delta_{G^v}(Y)|\qquad\text{for every } Y\subseteq V-v .
\end{equation}
Equivalently, by Gale's supply--demand theorem: in $G^v$, the unmarked even vertices other than $v$
can each send one unit along edge-disjoint paths to distinct marked vertices or to $v$, and $v$ may
receive any number of units.
\end{lemma}

\begin{proof}
$b_M(Y)=|Y\cap\Od|+2|Y\cap M|=|Y|-x_M(Y)$, directly from the definitions, since every vertex lies
in $\Od$ or $\Ev$. So (c) reads
$x_M(Y)\le |Y|-e(v,Y)+|\delta_{G-v}(Y)|=|\delta_{G^v}(Y)|$. For the flow statement, apply max-flow
min-cut to $G^v$ with these supplies, unit demands at the marked vertices, and an uncapacitated sink
at $v$. The cuts that do not separate $v$ from the sink are exactly the sets $Y\subseteq V-v$.
\end{proof}

A non-neighbour of $v$ is joined to $v$ in $G^v$. So the content of~\eqref{eq:signed} at $v$ is that
the unmarked even neighbours of $v$ can be routed in $G-v$. Each can go to a marked neighbour of $v$
(one slot each), to a marked non-neighbour (two slots) or to an odd non-neighbour (one slot). The
constant profile $b\equiv 1$, which is $x\equiv 0$, satisfies (c) trivially. So CS below is a parity
rounding problem: round $b\equiv1$ on $\Ev$ to values in $\{0,2\}$ while keeping $n$ cut systems
satisfied at once.

\begin{lemma}[parity]\label{lem:parity}
For $v\in V$ and $Y\subseteq V-v$, we have $|\delta_{G^v}(Y)|\equiv |Y\cap\Ev|\pmod 2$. In particular
both sides of~\eqref{eq:signed} have the same parity, so a violated inequality is violated by at
least two.
\end{lemma}

\begin{proof}
$|\delta_{G^v}(Y)|=|\delta(Y)|+|Y|-2e(v,Y)\equiv\sum_{y\in Y}(d(y)+1)$, and $d(y)+1$ is odd exactly
when $y\in\Ev$. Also $x_M(Y)\equiv|Y\cap\Ev|$.
\end{proof}

\section{Selection statements and matroid structure}\label{sec:consequences}

\begin{definition}
\textbf{(CS$_0$)}: $G$ has a cut-feasible half-set. \textbf{(CS)}: for every even vertex $r$, $G$ has a
cut-feasible half-set containing $r$. \textbf{(PS)}: $\Ev$ has a \emph{good pairing}, that is, a
partition into pairs such that every half-set meeting each pair in exactly one vertex is
cut-feasible. Clearly PS$\Rightarrow$CS$\Rightarrow$CS$_0$.
\end{definition}

\begin{remark}[relation to Gallai's conjecture]\label{rem:converse}
A realisable half-set gives a partition into $n/2$ paths. By Proposition~\ref{prop:nec},
cut-feasibility is necessary for a half-set to be realisable. So CS at $(G,r)$ follows from the
\emph{simple-profile rooted statement} (R$_s$): $r$ lies in a realisable half-set. Gallai's
conjecture is equivalent to the rooted test (R) (Lean \texttt{gallai\_iff\_rootedTest}, ES~\S1).
However, (R) gives a partition into $n/2$ paths with an end at $r$ that need not have a simple profile
(ES Corollary~6). Theorem~\ref{thm:ps} proves CS outright. Section~\ref{sec:hard} shows that a
cut-feasible half-set need not be realisable. So cut selection by itself says nothing about Gallai's
conjecture, and we make no such claim. Its role in the project is as a guide for an existential
descent towards (R$_s$) (\texttt{research/synthesis.md}~\S13).
\end{remark}

\subsection{Caps, matroids and EMC}\label{sec:caps}

\begin{conjecture}[EMC, endpoint matroid conjecture; \texttt{research/experiments.md} E6]\label{conj:emc}
If $G$ is critical, meaning $p(G)=n/2$, then the realisable half-sets are the bases of a loopless
matroid on $\Ev$.
\end{conjecture}

It holds for every critical graph of order at most $8$ (E15). Where Question~\ref{conj:emm} has a
positive answer (for instance at order $\le 8$), the realisable half-sets are the cut-feasible ones. The inequalities are upper bounds
$|M\cap X|\le f(X)$, and Edmonds' theorem turns a submodular cap into a matroid. So one asks whether
the cap is submodular. The answer has two parts.

\begin{theorem}[one vertex at a time]\label{prop:pervertex}
For each vertex $v$, the half-sets satisfying (c) at $v$ form the set $\cF_v$ of bases of a matroid
on $\Ev$.
\end{theorem}

\begin{proof}
For $Y\subseteq V-v$ put $\lambda_v(Y)=e(v,Y)-|\Od\cap Y|-|\delta_{G-v}(Y)|$. Condition (c) at $v$
reads $2|M\cap Y|\ge\lambda_v(Y)$. The function $\lambda_v$ is modular minus a cut function, hence
supermodular, and it is even. Evenness holds because
$|\delta_{G-v}(Y)|\equiv\sum_{y\in Y}(d(y)-[y\in N(v)])\equiv|Y\cap\Od|+e(v,Y)$. For
$X\subseteq\Ev-v$ put
\[
\ell_v(X)=\max\{\lambda_v(Y)/2:\ Y\subseteq V-v,\ Y\cap\Ev=X\}.
\]
This projection of a supermodular function is supermodular. If $Y_1$ and $Y_2$ attain $\ell_v(X_1)$
and $\ell_v(X_2)$, then $Y_1\cup Y_2$ and $Y_1\cap Y_2$ project to $X_1\cup X_2$ and $X_1\cap X_2$.
Condition (c) at $v$ is $|M\cap X|\ge\ell_v(X)$ for all $X\subseteq\Ev-v$. Since
$\lambda_v(Y)\le|Y|-|Y\cap\Od|=|X|$, we get $\ell_v(X)\le\lfloor |X|/2\rfloor$.

In terms of the unmarked set $N=\Ev\setminus M$, the condition says $|N\cap X|\le g_v(X)$ for all
$X\subseteq\Ev$, where $g_v(X):=|X|-\ell_v(X\setminus\{v\})$. (For $X\ni v$ this is the condition at
$X-v$ together with $|N\cap\{v\}|\le1$.) The function $g_v$ is submodular on $2^{\Ev}$, and
$g_v(X)\ge\lceil |X|/2\rceil\ge 0$. By Edmonds' theorem~\cite{Edm70}, the sets $N$ with
$|N\cap X|\le g_v(X)$ for all $X$ are the independent sets of a matroid $\mathcal I_v$. Its rank is
$\min_X(|\Ev\setminus X|+g_v(X))\ge\min_X(e-|X|+|X|/2)\ge e/2$. So the complements of the members of
$\cF_v$ are the independent sets of size $e/2$, which are the bases of the truncation of
$\mathcal I_v$ to rank $e/2$. Complements of bases are bases of the dual.
\end{proof}

The cut-feasible family is $\cF=\bigcap_v\cF_v$, an intersection of $n$ matroid basis families. In
the form~\eqref{eq:one} the caps are on the marks. For $v$ even, the inequalities at $v$ bound
$|M\cap X|$ for $X=(W\cup\{v\})\cap\Ev$; for $v$ odd they bound it for $X=W\cap\Ev$. Let $f_v(X)$ be
the least such bound, $\tfrac12\min_W(|\delta(W)|-|W\cap\Od|+2k-d(v)-[v\in\Od])$. Parity makes the
halving exact, so $f_v$ is again a projection of a submodular function and is submodular on its
domain. The global cap is $f=\min_v f_v$. The minimum of submodular functions need not be
submodular, and computation confirms that it fails here (Table~\ref{tab:caps}).
\begin{itemize}
\item The global cap $f$ is not submodular for $1{,}063$ of the $10{,}893$ connected graphs of order
eight. In $70$ of these, the cap system $\{I:|I\cap X|\le f(X)\}$ is not even a matroid.
\item $\cF$ fails to be a matroid basis family in $48$ graphs of order eight, and in one of order
six. All of them are non-critical, so this does not contradict EMC. An example is
\texttt{GsaAD\{}: two adjacent vertices $0,7$ of degree $6$, four common neighbours $2,3,4,5$ and the
path $0,1,6,7$. Here $X=\{1,4,5,6\}$ and $Y=\{0,1,4,7\}$ are in $\cF$, but neither
$X-5+0$ nor $X-5+7$ is.
\item For every one of the $1{,}082$ critical graphs of order eight, $\cF$ is a matroid basis family.
In $1{,}071$ of them $\cF=\cF_v$ for a single vertex $v$, so EMC follows there through
Theorem~\ref{prop:pervertex}. In $1{,}057$ the global cap is submodular. Four critical graphs are
explained by neither:
\texttt{Gs`ALo}, \texttt{GsQcaO}, \texttt{GqHTTS} and \texttt{GqhTTS}.
\end{itemize}
So \textbf{a positive answer to Question~\ref{conj:emm} would not formally imply EMC}. What it would
give is that EMC is equivalent, for
critical graphs, to the statement that $\bigcap_v\cF_v$ is a loopless matroid basis family.
Looplessness is CS.

\begin{question}\label{q:emc}
Why is $\bigcap_v\cF_v$ a matroid basis family when $G$ is critical? Is there always a set of
vertices whose conditions imply all others and whose caps combine submodularly?
\end{question}

\begin{remark*}[matroid intersection]
Every matroid $\mathcal I_v$ in the proof of Theorem~\ref{prop:pervertex} has
$g_v(X)\ge\lceil|X|/2\rceil$, so every set $A$ has rank at least $\lceil|A|/2\rceil$. For two vertices
$u,w$ and an even vertex $r$, Edmonds' matroid intersection theorem on $\Ev-r$ gives a common
independent set of size
$\min_A\bigl(r_u(A)+r_w(\Ev-r-A)\bigr)\ge\lceil|A|/2\rceil+\lceil(e-1-|A|)/2\rceil\ge e/2$. So the
conditions at any two vertices can be met together by a half-set containing $r$. This does not
extend to three matroids of this kind. On $\{1,2,3,4\}$, the partition matroids with blocks
$\{12,34\}$, $\{13,24\}$ and $\{14,23\}$ all have rank $\ge\lceil|A|/2\rceil$ on every $A$, yet they
have no common independent pair. At order $8$, cutting out $\cF$ needs up to six of the $\cF_v$: a
single one suffices in $10{,}798$ graphs and two in $75$. CS is nevertheless true for every graph
(Theorem~\ref{thm:ps}). In matroid terms, the theorem exhibits a partition matroid, that of a good
pairing, all of whose bases lie in every $\cF_v$.
\end{remark*}

\section{Cut selection}\label{sec:cs}

\begin{lemma}[good pairings]\label{lem:pairing}
A partition $P$ of $\Ev$ into pairs is good if and only if, for every $v$ and every $Y\subseteq V-v$,
\[
\sigma_P(Y)\le|\delta_{G^v}(Y)|,
\]
where $\sigma_P(Y)$ is the number of pairs with exactly one vertex in $Y$. Equivalently, the demand
graph $P$ satisfies the cut condition in $G^v$ for every vertex $v$.
\end{lemma}

\begin{proof}
For a transversal $M$ of $P$, a pair inside $Y$ contributes $1-1=0$ to $x_M(Y)$, a pair outside $Y$
contributes $0$, and a split pair contributes $\pm1$. So $x_M(Y)\le\sigma_P(Y)$, and
Lemma~\ref{lem:signed} gives sufficiency. Conversely, for given $v$ and $Y$, mark in each split pair
the vertex outside $Y$. Then $x_M(Y)=\sigma_P(Y)$. Every cut has a side avoiding $v$, so this is the
cut condition in $G^v$.
\end{proof}

\subsection{Cones}
The following answers the cone case.

\begin{lemma}[$T$-join pairing]\label{lem:tjoin}
In any graph $H$, the odd vertices can be partitioned into pairs joined by pairwise edge-disjoint paths
of $H$.
\end{lemma}

\begin{proof}
Induct on $|E(H)|$. If $H$ has a cycle, delete its edges. This changes no parity, so the induction
hypothesis applies. If $H$ is a forest with an edge, some tree component has two leaves $a$ and $c$.
Both are odd. Deleting the edges of the $a$--$c$ path makes $a$ and $c$ isolated and keeps all other
parities. Apply induction and add the pair $\{a,c\}$.
\end{proof}

\begin{theorem}[cones]\label{thm:cone}
Let $G$ have a universal vertex $u$ and even order $n$, and put $H=G-u$. A half-set $M$ is
cut-feasible if and only if $\Ev$ can be partitioned into pairs, each consisting of one vertex of $M$
and one of $\Ev\setminus M$, that are joined by pairwise edge-disjoint paths of $H$. Consequently,
every pairing of $\Ev$ by edge-disjoint paths of $H$ is good, such pairings exist, and PS and CS hold
for cones.
\end{theorem}

\begin{proof}
Here $d(u)=n-1$ is odd, so $u\notin\Ev$ and $\Ev(G)=\Od(H)$. In $G^u$ the vertex $u$ is isolated, so
by Lemma~\ref{lem:signed} the conditions at $u$ read $x_M(Y)\le|\delta_H(Y)|$ for $Y\subseteq V(H)$.
Since $x_M(V(H))=0$, applying them to $V(H)\setminus Y$ gives $|x_M(Y)|\le|\delta_H(Y)|$. This is the
Gale--Hoffman condition~\cite{Gale57,Hof60} for the supplies $x_M$ (unit supply at each unmarked even
vertex, unit demand at each marked one) to be routable by a flow of value at most one on each edge of
$H$. An integral such flow consists of edge-disjoint paths from unmarked to marked vertices, each such
vertex being the end of exactly one, together with cycles that can be discarded. Conversely, such
paths are a flow, and they imply the cut condition.

It remains to show that the conditions at $u$ imply those at every $v\neq u$. Use form (a). Write
$s(w)=2k-d(w)-b(w)$, so that \eqref{eq:one} at $(w,W)$ reads $b(W)-|\delta(W)|\le s(w)$. All
$s(w)\ge 0$, and $s(u)=2k-(2k-1)-1=0$. Let $W\subseteq V-v$. If $u\notin W$, then \eqref{eq:one} at
$(v,W)$ follows from \eqref{eq:one} at $(u,W)$, since $s(u)\le s(v)$. If $u\in W$, use form (b) with
$U=V-W$, which contains $v$ and lies in $V(H)$. Then $|\delta(U)|=|\delta_H(U)|+|U|$, and
$d(v)=1+d_H(v)\le 1+|\delta_H(U)|+(|U|-1)$. Also $b(v)\le b(U)$. Hence
$d(v)+b(v)\le b(U)+|\delta(U)|$.

Finally, any pairing by edge-disjoint paths of $H$ gives, for every choice of marked ends, the paths
required above. Lemma~\ref{lem:tjoin} applied to $H$ provides such a pairing, and choosing the end $r$
gives CS.
\end{proof}

So, for a cone, Question~\ref{conj:emm} asks whether \emph{a half-set is realisable if and only if
its marked and unmarked vertices pair up by edge-disjoint paths of $G-u$}. This is open; the
construction of Section~\ref{sec:hard} is not a cone.

\subsection{The guarded route lemma}

\begin{definition}
A \emph{route system} for a pairing $P$ of $\Ev$ consists of pairwise edge-disjoint paths $R_p$ of $G$
($p\in P$), with $R_p$ joining the two vertices of $p$. Let $J$ be the union of their edges. A vertex
is \emph{busy} if it lies on two or more routes. A \emph{passage} at $v$ is an unordered pair
$\{x,y\}$ such that $x,v,y$ are consecutive on a route. We write it $(x,y)$ when an orientation is
fixed inside a proof.

A passage $(x,y)$ at a busy vertex $v$ is \emph{guarded} if it has
\begin{itemize}
\item[(a)] a \emph{bypass}: a vertex $u\neq v$, which is not an end of another passage at $v$, with
$xu,yu\in E(G)\setminus J$; or
\item[(b)] \emph{private neighbours}: vertices $w_x,w_y\notin N(v)\cup\{v\}$ with
$xw_x,yw_y\in E(G)\setminus J$.
\end{itemize}
The private neighbours used by the passages at $v$ must be pairwise distinct; each end of each passage
has one. The route system is
\emph{guarded} if every passage at every busy vertex is guarded. It is a \emph{linkage} if the routes
are pairwise vertex-disjoint, in which case no vertex is busy.
\end{definition}

\begin{theorem}[guarded route lemma]\label{thm:guarded}
If a pairing $P$ of $\Ev$ has a guarded route system, then $P$ is good. In particular, every linkage
of $\Ev$ is a good pairing.
\end{theorem}

\begin{proof}
Fix $v$ and $Y\subseteq V-v$. By Lemma~\ref{lem:pairing} we must show
$\sigma_P(Y)\le|\delta_{G^v}(Y)|$. Both sides are congruent to $|Y\cap\Ev|$ modulo $2$: the right side
by Lemma~\ref{lem:parity}, and the left because $|Y\cap\Ev|=\sigma_P(Y)+2\cdot\#\{p\in P:p\subseteq
Y\}$. So it suffices to prove $\sigma_P(Y)\le|\delta_{G^v}(Y)|+1$.

Let $p=\{a,c\}$ be split, with $a\in Y$ and $c\notin Y$. Call $p$ \emph{good} if $R_p$ contains an
edge of $\delta_{G-v}(Y)$, and choose one such edge $\varepsilon_p\in J$. Distinct good pairs have
distinct $\varepsilon_p$, because routes are edge-disjoint.

Now suppose $p$ is bad. Follow $R_p$ from $a$. Every edge from $Y$ to $V-Y$ either goes to $v$ or lies
in $\delta_{G-v}(Y)$. So $R_p$ first leaves $Y$ along an edge $xv$ with $x\in Y$. If $c=v$, we call
$p$ a \emph{bad end}. There is at most one bad end, since $v$ lies in at most one pair. Otherwise $v$
is interior to $R_p$; let $y$ follow $v$. If $y\in Y$, the rest of $R_p$ would have to leave $Y$
again without returning to $v$, through an edge of $\delta_{G-v}(Y)$, which is impossible. So
$(x,y)$ is a passage at $v$ with $x\in Y$ and $y\notin Y\cup\{v\}$.

If $v$ lies on at most one route, at most one pair is bad. So
$\sigma_P(Y)\le|\delta_{G-v}(Y)|+1\le|\delta_{G^v}(Y)|+1$.

Let $v$ be busy. To each bad pair $p$ that is not the bad end we assign a \emph{token} $\tau_p$ from
$(\delta_{G-v}(Y)\setminus J)\cup(Y\setminus N(v))$, using the guard of its passage $(x,y)$:
\begin{itemize}
\item bypass $u$: take $\tau_p=xu$ if $u\notin Y$, and $\tau_p=yu$ if $u\in Y$;
\item private neighbours: take $\tau_p=xw_x$ if $w_x\notin Y$, and $\tau_p=w_x$ (a vertex of
$Y\setminus N(v)$) if $w_x\in Y$.
\end{itemize}
In each case the edge token lies in $\delta_{G-v}(Y)$ and not in $J$, since $u,w_x\ne v$.

We claim that tokens are distinct. Vertex tokens are distinct by hypothesis, and they are never
edges. Suppose an edge token $\tau_p=z_ph_p$ equals $\tau_{p'}=z_{p'}h_{p'}$ for $p\neq p'$. Here
$z_p\in\{x,y\}$ is the end of the passage of $p$ and $h_p$ is its bypass or private neighbour. We
cannot have $z_p=z_{p'}$, because then the edge $vz_p$ would lie on both routes. Hence $z_p=h_{p'}$
and $z_{p'}=h_p$. Now $z_{p'}\in N(v)$, so $h_p$ is not a private neighbour. So $h_p$ is a bypass,
and it is an end of the passage of $p'$ at $v$, which the definition forbids.

Therefore
$\sigma_P(Y)\le|\delta_{G-v}(Y)\cap J|+|\delta_{G-v}(Y)\setminus J|+|Y\setminus N(v)|+1
=|\delta_{G^v}(Y)|+1$.
\end{proof}

\begin{corollary}\label{cor:linkage}
PS (hence CS) holds for $G$ in each of the following cases.
\begin{enumerate}
\item $G$ has a path, not necessarily spanning, through all even vertices. Pair consecutive even
vertices along it. In particular, this covers every graph with a Hamiltonian path.
\item More generally, $G$ has vertex-disjoint paths covering $\Ev$, each containing an even number of
even vertices.
\item $G[\Ev]$ has a perfect matching. When $\Ev=V$, a linkage of $\Ev$ is exactly a perfect matching
of $G$.
\item $G$ is a cone, by Theorem~\ref{thm:cone}, or by Theorem~\ref{thm:guarded} with the apex as a
bypass of every passage.
\item $G=K_{2,m}$ with $m$ even. This has no linkage for $m\ge 4$. Let $p,q$ be the hubs and
$\ell_1,\dots,\ell_m$ the leaves. Pair $\{p,\ell_1\}$ and $\{q,\ell_2\}$ by edges, and pair
$\{\ell_{2i-1},\ell_{2i}\}$ ($i\ge 2$) through $p$. The only busy vertex is $p$, and $q$ bypasses each
of its passages.
\end{enumerate}
\end{corollary}

\subsection{Cut selection holds for every graph}

Fix a set $J$ of edges whose odd-degree vertices are exactly the even vertices of $G$ (a
\emph{$T$-join} for $T=\Ev$), and put $F=G-J$, the \emph{free} edges. For a vertex $v$ let $F^v$ be
the graph obtained from $F$ by deleting the edges at $v$ and joining $v$ to every vertex of
$V-v-N(v)$. For $Y\subseteq V-v$ we have
$|\delta_{F^v}(Y)|=|\delta_{F-v}(Y)|+|Y\setminus N(v)|$, and so
\begin{equation}\label{eq:split}
|\delta_{G^v}(Y)|=|\delta_{J-v}(Y)|+|\delta_{F^v}(Y)| ,
\end{equation}
where $\delta_{J-v}(Y)$ is the set of edges of $J$ between $Y$ and $V-Y-v$.

\begin{lemma}[free degrees]\label{lem:free}
Every vertex has odd degree in $F$. The odd vertices of $F^v$ other than $v$ are exactly the
$J$-neighbours of $v$, and $v$ itself is odd in $F^v$ if and only if $v\in\Ev$.
\end{lemma}

\begin{proof}
$d_J(w)\equiv[w\in\Ev]$ and $d(w)\equiv[w\notin\Ev]$, so $d_F(w)$ is odd. For $w\neq v$,
$\deg_{F^v}(w)=d_F(w)-[vw\in F]+[w\notin N(v)]$. This is $d_F(w)$, which is odd, when $vw\in J$. It is
$d_F(w)-1$ when $vw\in F$ and $d_F(w)+1$ when $w\notin N(v)$, both even. Finally
$\deg_{F^v}(v)=n-1-d(v)$, which is odd exactly when $d(v)$ is even.
\end{proof}

For a route system with union $J$ (which is a $T$-join for $T=\Ev$), the \emph{transitions} at $v$
are the passages at $v$. If $v\in\Ev$ there is also the \emph{end edge} $vz$ of the route ending at
$v$.

\begin{lemma}[linked transitions]\label{lem:linked}
Let a route system for a pairing $P$ have union $J$. Suppose that for every vertex $v$ there are pairwise edge-disjoint paths of $F^v$, one joining $x$
to $y$ for each passage $\{x,y\}$ at $v$, and, if $v\in\Ev$, one joining $z$ to $v$ for the end
edge $vz$. Then $P$ is good.
\end{lemma}

\begin{proof}
Fix $v$ and $Y\subseteq V-v$. As in the proof of Theorem~\ref{thm:guarded}, every split pair is either
good, and then charged injectively to an edge of $\delta_{J-v}(Y)$, or bad. A bad pair is the route
ending at $v$ with its end edge $vz$, $z\in Y$, or it has a passage $(x,y)$ at $v$ with $x\in Y$ and
$y\notin Y\cup\{v\}$. Charge a bad pair to an edge of $\delta_{F^v}(Y)$ on its linking path. The path
from $x\in Y$ to $y\notin Y$, or from $z\in Y$ to $v\notin Y$, crosses $\delta_{F^v}(Y)$. Distinct bad
pairs have distinct transitions at $v$, hence edge-disjoint linking paths and distinct charges. So
$\sigma_P(Y)\le|\delta_{J-v}(Y)|+|\delta_{F^v}(Y)|=|\delta_{G^v}(Y)|$ by~\eqref{eq:split}, and
Lemma~\ref{lem:pairing} applies.
\end{proof}

The guards of Theorem~\ref{thm:guarded} are special linking paths. A bypass gives $x,u,y$, and private
neighbours give $x,w_x,v,w_y,y$. (Theorem~\ref{thm:guarded} also spends the parity unit of
Lemma~\ref{lem:parity} on the end edge.) The decisive observation is that, once $J$ is fixed, the
hypothesis at $v$ involves only the transitions at $v$. When $J$ is a forest, these can be chosen
independently at every vertex.

\begin{theorem}[cut selection]\label{thm:ps}
Every connected graph $G$ of even order has a good pairing. Hence PS and CS hold for every connected
graph of even order. A good pairing can be found in polynomial time as follows.
\begin{enumerate}
\item Let $J$ be the $T$-join, $T=\Ev$, of a spanning tree of $G$, or any other $T$-join that is a
forest (for example a minimum one), and let $F=G-J$.
\item At each vertex $v$ with $d_J(v)>0$, pair the odd vertices of $F^v$ by edge-disjoint paths of
$F^v$ (Lemma~\ref{lem:tjoin}). By Lemma~\ref{lem:free} these are the $J$-neighbours of $v$, together
with $v$ when $v\in\Ev$. Declare $vx,vy$ consecutive whenever $\{x,y\}$ is a pair, and declare $vz$
an end edge when $z$ is paired with $v$.
\item Follow these transitions. The resulting paths pair $\Ev$, and their pairing is good.
\end{enumerate}
\end{theorem}

\begin{proof}
At each vertex the transitions pair up the $J$-edges, leaving exactly one unpaired edge at the
vertices of odd $J$-degree, that is, at $\Ev$. Following transitions therefore decomposes $J$ into
edge-disjoint trails and closed trails, and every trail ends at two vertices of $\Ev$, each the end of
exactly one trail. Since $J$ is a forest, there are no closed trails, and no trail repeats a vertex.
So the trails are edge-disjoint paths pairing $\Ev$. Their transitions at each $v$ are linked in $F^v$
by construction. At a vertex with $d_J(v)=0$ there is nothing to link, because then $v\in\Od$ (by the
parity of $d_J$) and $v$ has neither passages nor an end edge. So Lemma~\ref{lem:linked} shows that
the pairing is good.
\end{proof}

So the cut inequality can always be met, by an explicit half-set, and even by every transversal of
an explicit pairing. The script \texttt{scripts/construction.py} runs the construction with random
spanning trees and checks the result against the cut enumeration. With three trees per graph, it
checks all $99$ connected graphs of order $6$ with $e>0$ (of $112$ connected graphs) and all
$10{,}893$ with $e>0$ of order $8$ (of $11{,}117$). It also covers the order-$10$ sample, the
bottleneck families, all trees of orders $8$--$12$ and $150$ random graphs of order $12$. There are no
failures. An independent re-implementation of inequality~\eqref{eq:one} and of the construction
(audit for this manuscript) tested about $50{,}000$ further graphs of orders $8$--$12$ in twelve
families, and found no failure.

\begin{remark*}[certificates and guards]
Before Theorem~\ref{thm:ps} was found, \texttt{scripts/certificates.py} certified PS at every order
$\le 8$ by Hamiltonian paths ($9{,}848$ graphs at order $8$), other linkages ($996$) and guarded
route systems ($49$) (Table~\ref{tab:cert}). Whether every graph has a route system with the
restricted guards of Theorem~\ref{thm:guarded} is now immaterial. For instance, at order $8$ there are $31$ graphs in which no decomposition of a minimum $T$-join is
guarded. Yet the proof of Theorem~\ref{thm:ps} works for any forest $T$-join, including a minimum one,
so some decomposition always has linked transitions.
\end{remark*}

\subsection{Remarks}
\begin{enumerate}
\item \emph{The auxiliary networks.} An apex joined to the vertices of positive slack is realised
exactly, but separately for each $v$. In $G^v$ the vertex $v$ is joined to its $n-1-d(v)=s(v)+b(v)-1$
non-neighbours, and in the proof it is $F^v$, the free part of $G^v$, that carries the linking paths.
No single global network can work. The conditions at different vertices define different matroids
(Theorem~\ref{prop:pervertex}), and their intersection is not a matroid in general. Simultaneity is
achieved instead by fixing $J$ first: the $n$ cut systems then decouple into independent choices of
transitions.
\item \emph{Parity.} Lemma~\ref{lem:free} is the heart of the argument. Because $J$ is a $T$-join for
exactly the even vertices, every vertex keeps an odd number of free edges. The odd vertices of $F^v$
are then precisely the ends of the transitions at $v$, so the $T$-join lemma links them.
\end{enumerate}

\section{Cut-feasibility is not sufficient, and realisability is hard}\label{sec:hard}

The answer to Question~\ref{conj:emm} is positive at every order up to $8$ and in the dense range.
Both counterexamples below use \emph{saturated} vertices, that is, vertices with $b(v)=d(v)$. Every
edge at a saturated vertex ends a path there, so no path passes through it. For a simple profile
this needs $d(v)\le2$.

\begin{example*}[order $12$; missing-move Theorem~7.5]
Let $G=$\,\texttt{KsOeBW??G`?A}, with edges $01,02,03,06,14,16,17,25,27,47,57$ and the tail
$3\text{--}10\text{--}8\text{--}9\text{--}11$, and let $M=\{3,4,8,9,10\}$. Then $M$ is cut-feasible
(\texttt{tools/cutfeas.c}, \texttt{tools/cut\_test.py}) but not realisable (the exact oracle, and an
independent brute force). The vertices $3,4,8,9,10$ are saturated. So $3\text{--}10$, $10\text{--}8$,
$8\text{--}9$ and $9\text{--}11$ are one-edge paths, and the remaining eleven edges must form the
other two of the six paths: one from $3$ and one from $2$, both ending at $4$. The eight-vertex core
thus sits below its own budget, where a two-linkage obstruction prevents this and no cut detects it.
Twelve such graphs of order $12$ are known (\texttt{scripts/evidence.py counterexamples}).
\end{example*}

\subsection{A reduction from Hamiltonicity}
The construction and the observation that it encodes Hamiltonicity are due to the project lead. We
give complete proofs.

\begin{definition}\label{def:GH}
Let $H$ be a connected cubic graph, and let $L(H)$ be its line graph: $4$-regular, with vertex set
$E(H)$. Fix a vertex $w^*$ of $L(H)$ with neighbours $u_1,\dots,u_4$. Put $R=V(L(H))\setminus\{w^*\}$
and $N=|R|=|E(H)|-1\ge5$. The graph $G(H)$ has vertex set $R\cup\{l_1,\dots,l_N\}$ and edge set
consisting of
\begin{itemize}
\item the edges of $L(H)-w^*$;
\item the path $l_1l_5l_6\cdots l_Nl_2$ and the edge $l_3l_4$;
\item the \emph{cross edges} $l_iu_i$ for $1\le i\le 4$.
\end{itemize}
Vertices of $R$ have degree $4$ and the $l_i$ have degree $2$. So $n=2N$, $\Ev=V$, and the
\emph{gadget} $\mathcal L=\{l_1,\dots,l_N\}$ is a half-set. Its profile $b$ is $2$ on $\mathcal L$ and
$0$ on $R$, and it asks for $N$ paths.
\end{definition}

\begin{lemma}[cuts of line graphs of cubic graphs]\label{lem:linecuts}
Let $H$ be a connected cubic graph and $E(H)=A\sqcup B$ with $A,B\ne\emptyset$. Then
$|\delta_{L(H)}(A)|=2m(A)$, where $m(A)$ is the number of vertices of $H$ incident with edges of both
$A$ and $B$. Consequently $L(H)$ is $4$-edge-connected if and only if $H$ is bridgeless.
\end{lemma}

\begin{proof}
Two edges of the simple graph $H$ are adjacent in $L(H)$ exactly when they share one endpoint. At a
vertex whose three edges split $1+2$ between $A$ and $B$, exactly two adjacent pairs cross. At a
vertex with all three edges on one side, none do. This proves the formula, and $m(A)\ge1$ because $H$
is connected.

Suppose $m(A)=1$, with mixed vertex $x$. Let $e=xy$ be the edge at $x$ that is alone on its side,
say $e\in A$, so $x$'s other two edges lie in $B$. We show that $e$ is a bridge. Every vertex other
than $x$ has all its edges on one side. Take a path $y=p_0,p_1,\dots,p_t=x$ in $H-e$. The edge
$p_0p_1$ is in $A$, because $y\ne x$ is incident with $e\in A$. Each internal vertex $p_i\ne x$
passes the side on, so every edge of the path is in $A$. But the last edge ends at $x$, whose only
$A$-edge is $e$, which is not on the path. This is a contradiction, so $e$ is a bridge.

Conversely, let $xy$ be a bridge, and let $K$ be the component of $H-xy$ containing $x$. Take $A=E(K)$,
which is nonempty because $x$ has degree $2$ in $K$. The only mixed vertex is $x$, so
$|\delta_{L(H)}(A)|=2$. Finally, every edge cut of $L(H)$ is $\delta(A)$ for some such partition, so
$L(H)$ is $4$-edge-connected exactly when $m(A)\ge2$ always.
\end{proof}

\begin{proposition}[cut-feasibility]\label{prop:Lfeasible}
The gadget $\mathcal L$ is a cut-feasible half-set of $G(H)$ if and only if $L(H)$ is
$4$-edge-connected, that is, if and only if $H$ is bridgeless.
\end{proposition}

\begin{proof}
We use the arm form (c) of Lemma~\ref{lem:forms} at every vertex $v$ and every $Y\subseteq V-v$. Put
$f_v(Y)=b(Y)+|\delta_{G-v}(Y)|-e(v,Y)$; we need $f_v\ge0$.

\emph{$v\in R$.} Adding a gadget vertex $l$ to $Y$ never decreases $f_v$. Indeed, $b$ grows by $2$.
The cut $\delta_{G-v}$ loses at most $\deg_{G-v}(l)=2-[l\sim v]$ edges. And $e(v,\cdot)$ grows by
$[l\sim v]$. So the net change is at least $2-(2-[l\sim v])-[l\sim v]=0$. It therefore suffices to
take $Y\subseteq R-v$, where $b(Y)=0$. For such $Y$, the cross edges from $Y\cap\{u_1,\dots,u_4\}$
play the role of the edges to $w^*$ in $L(H)$. So $|\delta_{G-v}(Y)|=|\delta_{L(H)-v}(Y)|$ and
$e(v,Y)=e_{L(H)}(v,Y)$. With $X=Y\cup\{v\}$ and $d_{L(H)}(v)=4$,
$|\delta_{L(H)}(X)|=|\delta_{L(H)-v}(Y)|+4-e_{L(H)}(v,Y)$. Hence (c) at $v$ is equivalent to
$|\delta_{L(H)}(X)|\ge4$ for every $X$ with $v\in X\subseteq R$. As $v$ ranges over $R$, these $X$
are all the nonempty vertex sets of $L(H)$ avoiding $w^*$, and every edge cut of $L(H)$ has such a
side. So the conditions at the vertices of $R$ hold if and only if $L(H)$ is $4$-edge-connected.

\emph{$v\in\mathcal L$, assuming $L(H)$ is $4$-edge-connected.} Here $d(v)=2$. If $Y$ meets
$\mathcal L$, then $b(Y)\ge2\ge e(v,Y)$. If $Y\subseteq R$, then $e(v,Y)\le1$, with equality only when
$v=l_i$ and $u_i\in Y$. In that case $\delta_{G-v}(Y)\ne\emptyset$. Indeed, $L(H)-w^*$ is connected,
since a cut vertex $w^*$ of the $4$-regular $L(H)$ would leave a side with at most $3$ edges to the
rest. So $Y\subsetneq R$ has an edge to $R\setminus Y$. And $Y=R$ has the three cross edges to the
other $l_j$.

Finally, Lemma~\ref{lem:linecuts} translates $4$-edge-connectivity of $L(H)$ into bridgelessness of
$H$.
\end{proof}

\begin{proposition}[realisations]\label{prop:Lreal}
The half-set $\mathcal L$ is realisable in $G(H)$ if and only if $L(H)$ has a Hamiltonian
decomposition, that is, two edge-disjoint Hamiltonian cycles covering $E(L(H))$.
\end{proposition}

\begin{proof}
Let $\mathcal P$ realise $b$ with $N$ paths. Every gadget vertex has $b=d=2$, so every edge at it ends
a path there. An edge between two gadget vertices therefore ends paths at both of its ends, so it is a
path by itself. These are the $N-3$ edges of the path $l_1l_5\cdots l_Nl_2$ and the edge $l_3l_4$,
which gives $N-2$ paths. A cross edge $l_iu_i$ ends a path at $l_i$. The vertices of $R$ have $b=0$
and degree $4$, so each lies, as an interior vertex, on exactly two paths. The remaining two paths
$Q_1,Q_2$ cover $E(L(H)-w^*)$ and the four cross edges, and their four ends are $l_1,\dots,l_4$. Each
vertex of $R$ is visited twice in total and at most once by each path, so by both. So $Q_1-\{l_a,l_b\}$
is a Hamiltonian path of $L(H)-w^*$ from $u_a$ to $u_b$, and $Q_2-\{l_c,l_d\}$ is one from $u_c$ to
$u_d$, with $\{a,b,c,d\}=\{1,2,3,4\}$. Closing them through $w^*$ gives two edge-disjoint Hamiltonian
cycles of $L(H)$ that cover all its edges.

Conversely, in a Hamiltonian decomposition $C_1,C_2$ of $L(H)$, each cycle uses two of the four edges
at $w^*$. Deleting $w^*$ and extending through the cross edges gives two paths $l_a\cdots l_b$ and
$l_c\cdots l_d$. Together with the $N-2$ one-edge gadget paths, they realise $b$.
\end{proof}

\begin{theorem}[Kotzig~\cite{Kot57}; see also~\cite{Ber78}]\label{thm:kotzig}
A cubic graph $H$ is Hamiltonian if and only if its line graph $L(H)$ has a Hamiltonian decomposition.
\end{theorem}

\begin{theorem}[hardness]\label{thm:hard}
Let $H$ be a connected bridgeless cubic graph. Then the gadget $\mathcal L$ is a cut-feasible half-set
of $G(H)$, and it is realisable if and only if $H$ is Hamiltonian. Consequently, deciding whether a
given half-set of a given connected graph is realisable is NP-complete. This remains true for graphs
with all degrees in $\{2,4\}$ and half-sets that satisfy every cut inequality~\eqref{eq:one}. Unless
P$\,=\,$NP, no condition checkable in polynomial time characterises the realisable half-sets.
\end{theorem}

\begin{proof}
Cut-feasibility is Proposition~\ref{prop:Lfeasible}. Realisability is Proposition~\ref{prop:Lreal}
together with Theorem~\ref{thm:kotzig}. Membership in NP holds because a realisation is a certificate.
Hardness follows because $G(H)$ has size $O(|V(H)|)$, and deciding Hamiltonicity is NP-complete
already for $3$-connected cubic planar graphs~\cite{GJT76}, which are bridgeless. Cut-feasibility, in
contrast, is decidable in polynomial time by the max-flow form (d) of Lemma~\ref{lem:forms}.
\end{proof}

\begin{example*}[order $28$]
For the Petersen graph $P$ (bridgeless, cubic, not Hamiltonian), $G(P)$ has $n=28$ and $m=42$, every
vertex is even, and every slack $2k-d(v)-b(v)$ equals $24$. The gadget is cut-feasible but not
realisable. \texttt{scripts/petersen\_counterexample.py} checks this directly: $L(P)$ has $160$
Hamiltonian cycles and no Hamiltonian decomposition, and no risky $W$ violates~\eqref{eq:one}.
\texttt{scripts/kotzig.py} checks Proposition~\ref{prop:Lfeasible} by max-flow. It checks
Proposition~\ref{prop:Lreal} with the exact oracle where $n\le16$ ($K_4$, the prism and $K_{3,3}$), and
Theorem~\ref{thm:kotzig} by search where $L(H)$ has at most $24$ vertices. The graphs
used are $K_4$, the prism, $K_{3,3}$, the cube, the Petersen graph, six random cubic graphs of orders
$8$--$18$, and a cubic graph with a bridge, for which $\mathcal L$ is not cut-feasible.
\end{example*}

\begin{remark}[what remains]\label{rem:what}
The counterexamples use saturated vertices of degree $2$. Whether cut-feasibility suffices at minimum
degree at least $3$ (\emph{EMM$_3$}) is open. The evidence is thin: at order $\le8$, every half-set
of a graph with $\delta\ge3$ is realisable, so the inequality never binds there
(Section~\ref{sec:evidence}). Replacing the gadget by degree-$3$ forcing structures would very likely
carry Theorem~\ref{thm:hard} to $\delta\ge3$ (\texttt{research/synthesis.md}~\S12). Theorem~\ref{thm:hard}
concerns a \emph{prescribed} half-set, and it does not bear on the existence of some realisable
half-set. In $G(H)$, for instance, $\mathcal L$ starves $R$ of ends: fourteen even vertices and only
four boundary edges in the Petersen case. Marks placed inside $R$ remove the Hamiltonicity requirement
at once.
\end{remark}

\section{The dictionary}\label{sec:dictionary}

In each item below, the first sentence gives the source statement, as extracted from the PDFs. The
rest is our reading in terms of~\eqref{eq:one}.

\begin{enumerate}
\item \textbf{SP Theorem 14.1 (dense case: cuts never bind).} For $\varepsilon>0$ and $r$, and $n$
large: if $\delta(G)\ge(1/2+\varepsilon)n$ and $b\le r$ has $b(V)=2k$, then $b$ is realised by $k$
paths if and only if the local conditions hold. The following proposition shows that in this range
inequality~\eqref{eq:one} is implied by the local conditions. So in the dense range the answer to
Question~\ref{conj:emm} is positive, for all bounded profiles and not only simple ones.

\begin{proposition}\label{prop:dense}
If $b\le r$ satisfies the local conditions and $\delta(G)\ge(n+r-1)/2$, then (c) holds at every
vertex. In particular, if $\delta(G)\ge(n+1)/2$, every half-set is cut-feasible.
\end{proposition}

\begin{proof}
Let $Y\subseteq V-v$. If $|Y|\le\delta-1$, each $y\in Y$ has at least $d(y)-1-(|Y|-1)\ge 1$
neighbours in $V-Y-v$. So $|\delta_{G-v}(Y)|\ge|Y|\ge e(v,Y)$. If $|Y|\ge\delta$, put $Z=V-v-Y$, so
$|Z|\le n-1-\delta$. Then (c) is equivalent to
$b(Z)-|\delta_{G-v}(Z)|\le 2k-d(v)-b(v)+e(v,Z)$, whose right side is non-negative. Each $z\in Z$ has
at least $\delta-|Z|\ge 2\delta-n+1\ge r$ neighbours in $V-Z-v$, so $|\delta_{G-v}(Z)|\ge r|Z|\ge
b(Z)$.
\end{proof}

SP's sharpness example consists of two cliques of order $2h+1$ joined by one edge, so $\delta=n/2-1$.
It puts one end at each end of the bridge, two at the other vertices of one clique, and none at the
other vertices of the second clique $B$. This violates~\eqref{eq:one} at every unmarked vertex $v$ of $B$. In arm form, take $Y=B-v$: then
$e(v,Y)=2h$ exceeds $b(Y)+|\delta_{G-v}(Y)|=1+1$ once $h\ge2$. So the two regimes are separated exactly
by the cut inequality.

\item \textbf{ED Proposition 45 and Theorem 46 (necessity and sharpness).} Proposition 45: if $G$ is
$\kappa$-connected with $\kappa\ge2$ and at most $2\kappa$ even vertices, then every choice of marks
meets every inequality~\eqref{eq:one}. Whether every half-set of such a graph is realisable is open;
by Theorem~\ref{thm:hard} it cannot follow from the cut inequality alone. Theorem 46: for $s\ge1$, there are graphs of connectivity $s+1$ and large
minimum degree with $2h$ even vertices in which prescribing two ends at $h$ specified even vertices
costs $F\ge k+(h-s)/2$ paths. The lower bound is \eqref{eq:one} with $k+j$ paths in place of $k$, at
$W=$ one shore and $v=$ the apex.

\item \textbf{TVS Theorems 1 and 7 (exact gluing at a 2-separator $\{u,v\}$).} Every attainable global
endpoint vector arises from two local records, three-piece joins $i,j$, and a matching of residual
path types in a $5\times5$ table. The pole counts are $b_A(z)+b_B(z)-2(i+j+m_z)$. The only global
ingredient is the matching number, which the paper computes by its K\"onig dual
$\min_I(\sum_{i\notin I}a_i+\sum_{j\in N(I)}b_j)$, a min-cut. So TVS's exact gluing is a cut
computation at a two-vertex separator, while~\eqref{eq:one} gives only necessary conditions. The
graphs of Section~\ref{sec:hard} show that the gap between the two can be NP-hard to close in
general.

\item \textbf{RE Theorem 13.1 (simultaneous rotation criterion).} Fix a partition $P$, a vertex $z$
and a set $A$ of non-neighbours of $z$. Spokes $zi$ ($i\in I$) can be restored by Fan's rotations and
appends if and only if a bipartite graph between $A$ and the rows omitting $z$ has a matching covering
$I$. This is Hall's condition, and it is the arm condition (d) at $z$, restricted to the arms that
rotations can build from the fixed partition $P$. Each new spoke needs its own arm to a row end not
at $z$, and the trees $T_i$ are the reachable arm ends.

\item \textbf{MR Corollary 39 (vertex insertion).} A new vertex $x$ is joined to $d$ ends of an
$s$-path partition. Let $r$ be the largest number of these sites that can be assigned to distinct old
paths ending there, and put $q=d-r$. Then $s+\lceil q/7\rceil$ paths suffice, with at least
$d-2q+2\lceil q/7\rceil$ ends at $x$. The denominator becomes $8$ by the remark after MR Theorem 175.
Here $q$ is the Hall deficiency of the trivial arm system at $x$ (arms of length $0$ into old ends).
Corollary 39 therefore says that a deficiency of $q$ in the arm condition at the inserted vertex,
relative to the old partition, costs at most $\lceil q/7\rceil$ extra paths.

\item \textbf{JM-repair Proposition 6.4 (Hoffman cuts).} The two-vertex extension step (Theorem~B) is
reduced to an admissible record. That in turn is reduced, through Hoffman's circulation
theorem~\cite{Hof60}, to two families of cut inequalities (UH), (LH) indexed by $B\subseteq V$. This is
the archive's one existing exact min-max for an endpoint-slot problem: slots at vertices must be
matched to colours, and feasibility is decided by cuts. Question~\ref{conj:emm} asked for the analogous
statement with paths in place of colour classes. Theorem~\ref{thm:hard} shows that for paths no such
statement can hold.

\item \textbf{This project.} The hub, flower and last-edge results of ES \S7 have the same
supply-versus-cut shape at a single vertex. They compare the ends available near a vertex with the
edges that must leave a small set around it.
\end{enumerate}

\section{Related work}\label{sec:related}

\textbf{Gallai's conjecture.} Lov\'asz~\cite{L68} proved that every graph on $n$ vertices decomposes
into at most $\lfloor n/2\rfloor$ paths and cycles. When all degrees are odd, the decomposition
consists of $n/2$ paths, which is the case $e=0$ of Question~\ref{conj:emm}. Pyber~\cite{Pyb96} proved Gallai's
conjecture when the even vertices induce a forest, and Fan~\cite{Fan05} extended this to wider
classes. Bonamy and Perrett~\cite{BP19} proved it for maximum degree at most $5$, and Botler and
Sambinelli~\cite{BS21} survey the area. These results bound the number of paths without prescribing
the ends. As far as two literature searches show (\texttt{research/synthesis.md}~\S7), no
characterisation of path partitions with prescribed end counts by cut conditions has appeared.
Theorem~\ref{thm:hard} explains why none can. Corollary~\ref{cor:linkage}(3) ($G[\Ev]$ has a perfect matching) is reminiscent
of Pyber's and Fan's hypotheses on the subgraph induced by the even vertices. Here, however, it
yields only the cut conditions, not a realisation.

\textbf{The Lov\'asz--Cherkassky theorem.} In an inner-Eulerian graph with terminal set $T$, the
maximum number of edge-disjoint $T$-paths is $\frac12\sum_{t\in T}\lambda(t,T-t)$
(Lov\'asz~\cite{L76}, Cherkassky~\cite{Ch77}; Frank~\cite{FrankBook} gives a splitting-off proof).
It is the classical min-max theorem closest to Question~\ref{conj:emm}, in two respects.
\begin{itemize}
\item \emph{The upper bound is charged to cuts in the same way.} A $T$-path leaving $t$ crosses every
cut separating $t$ from $T-t$. In the proof of Proposition~\ref{prop:nec}, each arm leaving $v$
either ends in $U$ or crosses $\delta(U)$.
\item \emph{The content is simultaneity.} Each $\lambda(t,T-t)$ is a single Menger number, and the
theorem says that all of them are attained, up to the factor $\frac12$, by one packing. In the same
way, each vertex condition here is a single Menger statement, $\lambda_{G_b}(v,t)=\deg_{G_b}(v)$
(Remark~\ref{rem:euler}). Question~\ref{conj:emm} asked whether all of them are attained by one
partition. Theorem~\ref{thm:ps} shows that they can all be met by one half-set, and
Theorem~\ref{thm:hard} that this does not suffice for a partition.
\end{itemize}
The analogy concerns the proof structure; it is not a reduction. Attaching $b(z)$ pendant terminals at
each vertex $z$ gives an inner-Eulerian graph for a simple profile. There, a packing of $k$
edge-disjoint $T$-paths that uses every end slot once is easy to find: pair the odd vertices by a
$T$-join (Lemma~\ref{lem:tjoin}), and pair the two slots of each marked vertex with each other. The
substance of realisability is that the packing can be chosen to \emph{cover every edge} with
non-trivial paths, so that no even subgraph is left over. In $G(H)$ this is a Hamiltonian
decomposition problem.

\textbf{Splitting off.} Mader's theorem~\cite{Mad78} is as follows. Let $s$ be a vertex of degree
other than $3$ that is not incident with a cut edge. Then some pair of edges at $s$ can be split off
(replacing $xs,sy$ by $xy$) without decreasing $\lambda(x,y)$ for any $x,y\neq s$. Frank~\cite{Fr92}
gave a short proof. The global version, for $k$-edge-connectivity, is due to Lov\'asz
(see~\cite{FrankBook}). In our setting this gives one clean fact.

\begin{proposition}\label{prop:mader}
Let $b$ be a simple profile and $v\in V$. Completely split off the vertex $v$ in $G_b$, deleting any
loops that arise, so that the $\lambda$-values in the resulting multigraph $G'_b$ are preserved. If
$b$ satisfies the conditions (b) at every vertex $w\neq v$, then $G'_b$ still satisfies
$\lambda(w,t)=\deg(w)$ for every $w\neq v$.
\end{proposition}

\begin{proof}
$G_b$ is Eulerian, so it has no cut edge. Every vertex has even degree, so it is never $3$. Splitting
off a pair at $v$ keeps the graph Eulerian, and deleting a loop does not change any $\lambda$. So
Mader's theorem applies repeatedly until $v$ is isolated. Splitting $xv,vy$ into $xy$ keeps the
degrees of $x$ and $y$. A loop arises only when $x=y$, that is, from two parallel edges at $v$. Since
$G$ is simple and split edges avoid $v$, these can only be the two edges $vt$ of a marked $v$. So the
only degrees that change are those of $v$ and $t$. Hence $\lambda'(w,t)\ge\lambda(w,t)=\deg(w)=\deg'(w)
\ge\lambda'(w,t)$ for every $w\neq v$.
\end{proof}

So the cut conditions survive the elimination of any vertex. The slack $2k-d(v)-b(v)$ of the anchor
$v$ is twice the number of paths that miss $v$, and those paths are untouched by the splitting. What splitting loses is the constraint that a path passes
through $v$ at most once. In $G'_b$, the cycles through $t$ must use at most one split edge from $v$
each. Any splitting-off argument for realisability would therefore have to choose Mader pairs
\emph{compatibly} with the cycle structure, and Theorem~\ref{thm:hard} shows that in general this
cannot be done by cut considerations. For CS, the
analogous step would be to split at $t$ in $G+t$ with $t$ joined once to every vertex (the
half-integral point $b\equiv1$). A good pairing is a complete splitting of the $t$-edges at $\Ev$
whose pairs satisfy the cut condition in every $G^v$ (Lemma~\ref{lem:pairing}). We have not found a
splitting theorem that guarantees this for all $v$ at once.

\textbf{Flows and $T$-joins.} Lemma~\ref{lem:forms} and Lemma~\ref{lem:signed} are instances of
max-flow min-cut and of Gale's and Hoffman's feasibility theorems~\cite{Gale57,Hof60}.
Lemma~\ref{lem:tjoin} is the path decomposition of a $T$-join~\cite{EJ73}, and
Theorem~\ref{prop:pervertex} is Edmonds' theorem on matroids induced by submodular
functions~\cite{Edm70}. Lemma~\ref{lem:pairing} says that a good pairing is a demand graph meeting
the cut condition in each $G^v$. When the cut condition suffices for edge-disjoint routing is a
classical question; Okamura and Seymour~\cite{OS81} settle it for planar supply graphs with all
terminals on one face. The guarded route lemma takes the opposite direction. It replaces routability
in every $G^v$ by a single routing in $G$ with local guards.

\section{Evidence}\label{sec:evidence}

All numbers are produced by the scripts in \texttt{scripts/}. They use the exact realisability oracle
\texttt{tools/realis.c} and the cut-feasibility enumerator \texttt{tools/cutfeas.c}, which is
cross-checked against \texttt{tools/cut\_test.py} at order six. The order-$10$ sample consists of $250$
connected graphs $G(10,p)$, $p\in\{.3,\dots,.8\}$, with seed $20261127$. The bottleneck families are
cliques joined by $1$--$3$ edges, theta graphs, cones over unions of paths, the eight-vertex example of
ED \S12.4, and 2-sums of small dense graphs. They have orders $6$--$10$, $82$ graphs after isomorphism
reduction.

\begin{table}[h]
\caption{Realisability versus cut-feasibility (\texttt{scripts/evidence.py}). A failure is a
cut-feasible half-set that is not realisable. ``Unsound'' counts realisable half-sets that violate
\eqref{eq:one}, and it must be $0$.}\label{tab:emm}
\begin{tabular}{lrrrr}
\toprule
Graphs & count & half-sets (non-realisable) & failures & unsound\\
\midrule
all connected, order 6 & 99 & 518 (18) & 0 & 0\\
\quad of which $\delta\ge3$ & 13 & 60 (0) & 0 & 0\\
all connected, order 8 & 10{,}893 & 101{,}714 (2{,}100) & 0 & 0\\
\quad of which $\delta\ge3$ & 2{,}475 & 19{,}422 (0) & 0 & 0\\
random, order 10 & 250 & 4{,}222 (8) & 0 & 0\\
\quad of which $\delta\ge3$ & 149 & 2{,}590 (0) & 0 & 0\\
bottleneck families, orders 6--10 & 82 & 1{,}568 (208) & 0 & 0\\
\quad of which $\delta\ge3$ & 22 & 324 (62) & 0 & 0\\
two blobs, order 10, $\delta\ge3$ & 60 & 1{,}218 (220) & 0 & 0\\
counterexamples of missing-move Thm 7.5, order 12 & 12 & 1{,}932 (269) & 12 & 0\\
\bottomrule
\end{tabular}
\end{table}

\textbf{Minimum degree at least three.} At orders $6$ and $8$, and in the order-$10$ random sample,
\emph{every} half-set of a graph of minimum degree at least $3$ is realisable, so the cut inequality
never excludes anything there. All $2{,}100$ non-realisable half-sets at order $8$ live in graphs with
a vertex of degree at most $2$. The rows in which the inequality does exclude half-sets at minimum
degree $\ge3$ are the bottleneck families and the two-blob graphs (two random dense blobs joined by
$1$--$3$ edges). In those rows every half-set it admits is realisable. The order-$28$ example of
Section~\ref{sec:hard} and the twelve order-$12$ examples all have vertices of degree at most $2$.

\begin{table}[h]
\caption{CS and PS by enumeration (\texttt{scripts/cut\_selection.py}). At order $6$ the script also
checks Lemma~\ref{lem:forms}(c) and Lemma~\ref{lem:pairing} against the enumeration on every half-set
and pairing.}\label{tab:cs}
\begin{tabular}{lrrrr}
\toprule
Graphs & count & CS failures & PS failures & a good pairing of edges of $G$\\
\midrule
all connected, order 6 & 99 & 0 & 0 & 63\\
all connected, order 8 & 10{,}893 & 0 & 0 & 6{,}773\\
random, order 10 & 250 & 0 & 0 & 181\\
bottleneck families & 82 & 0 & 0 & 35\\
\bottomrule
\end{tabular}
\end{table}

\begin{table}[h]
\caption{Certificates for Theorem~\ref{thm:guarded} (\texttt{scripts/certificates.py}).}\label{tab:cert}
\begin{tabular}{lrrrrr}
\toprule
Graphs & count & Hamiltonian path & other linkage & guarded routes & none\\
\midrule
all connected, order 6 & 99 & 82 & 16 & 1 & 0\\
all connected, order 8 & 10{,}893 & 9{,}848 & 996 & 49 & 0\\
random, order 10 & 250 & 236 & 14 & 0 & 0\\
bottleneck families & 82 & 69 & 13 & 0 & 0\\
\bottomrule
\end{tabular}
\end{table}

The certificate search was also run outside the scripts' fixed families. It found certificates for
all trees of orders $8$, $10$ and $12$: $20$, $99$ and $538$ trees with $e\ge 2$, of which $1$, $8$ and
$62$ need guards. It also found them for $300$ sparse random graphs ($p\in[.15,.3]$) of each of the
orders $10$ and $12$, and for $K_{2,m}$ and $K_{2,m}$ plus the hub edge, $m=6,8,10$.

Every certified pairing in Table~\ref{tab:cert} and every tree is also checked against the enumeration
of \texttt{tools/cutfeas.c}, and all of them are good. As a direct test of Theorem~\ref{thm:guarded},
\texttt{certificates.py stress} draws random pairings with random edge-disjoint route systems in $400$
random graphs of orders $8$ and $10$. All $3{,}310$ guarded systems with a busy vertex and all
$1{,}691$ linkages give good pairings. Of the $7{,}227$ unguarded systems, $7{,}079$ still give good
pairings. So the guards are far from necessary; Theorem~\ref{thm:ps} explains why.

\begin{table}[h]
\caption{Caps (\texttt{scripts/caps.py}); order $8$. The columns are: global cap submodular; cap system a
matroid; $\cF$ a matroid basis family; $\cF=\cF_v$ for one $v$; critical. At order $6$: $98$ of $99$
families are matroid basis families, and all $34$ critical graphs are.}\label{tab:caps}
\begin{tabular}{cccccr}
\toprule
$f$ submod. & cap matroid & $\cF$ matroid & single $v$ & critical & graphs\\
\midrule
no & no & no & no & no & 48\\
no & no & yes & no & no & 20\\
no & no & yes & no & yes & 2\\
no & yes & yes & no & no & 16\\
no & yes & yes & no & yes & 2\\
no & yes & yes & yes & no & 954\\
no & yes & yes & yes & yes & 21\\
yes & yes & yes & no & yes & 7\\
yes & yes & yes & yes & no & 8{,}773\\
yes & yes & yes & yes & yes & 1{,}050\\
\bottomrule
\end{tabular}
\end{table}

The script also verifies Theorem~\ref{prop:pervertex} by brute force: every $\cF_v$ is a matroid
basis family, at orders $6$ and $8$.

\section{Formalisation}
The shared Lean definitions are fixed in \texttt{research/connections.md}~\S U. We recommend the arm
form (c) as the definition, since its quantifiers are the simplest, with (a) proved equivalent.
Lemma~\ref{lem:parity}, Lemma~\ref{lem:pairing}, Lemma~\ref{lem:free}, Lemma~\ref{lem:linked} and
Theorem~\ref{thm:ps} are finite and elementary, so Theorem~\ref{thm:ps} (\texttt{cs\_all}) is the
natural target. Lemma~\ref{lem:linecuts} and Propositions~\ref{prop:Lfeasible} and~\ref{prop:Lreal}
are also elementary. Kotzig's theorem would be an imported hypothesis.

\section{Open problems}
\begin{enumerate}
\item EMM$_3$: is every cut-feasible half-set of a connected graph of minimum degree at least $3$
realisable? Or does Theorem~\ref{thm:hard} extend to minimum degree $3$ via forcing gadgets?
\item Existence instead of characterisation. Theorem~\ref{thm:hard} rules out characterising the
realisable half-sets, but not proving that \emph{some} half-set containing a given $r$ is realisable,
which is (R$_s$). Can the good pairing of Theorem~\ref{thm:ps} guide a descent from $\Ev$, which is
realisable (ES Lemma~2), to one of its transversals, keeping every intermediate set realisable?
(\texttt{research/synthesis.md}~\S13.)
\item Question~\ref{q:emc}: why is the cut-feasible family a matroid for critical graphs at order
$\le 8$, and does this persist?
\item Which cut-feasible half-sets are transversals of good pairings? Their union is closed under
complementation, but $\cF$ need not be: it is not in $593$ graphs of order $8$. At order $\le 8$, the
union equals the complement-closed part $\{M\in\cF:\Ev\setminus M\in\cF\}$ except in two graphs of
order $8$. In cones it is all of $\cF$ (Theorem~\ref{thm:cone}).
\item (R)$\Rightarrow$(R$_s$) (Remark~\ref{rem:converse}): can a partition with an end at $r$ always be
made simple?
\end{enumerate}

\begin{thebibliography}{99}
\bibitem{Ber78} J.-C.~Bermond, Hamiltonian decompositions of graphs, directed graphs and hypergraphs,
Ann.\ Discrete Math.\ 3 (1978) 21--28.
\bibitem{BP19} M.~Bonamy and T.~J.~Perrett, Gallai's path decomposition conjecture for graphs of small
maximum degree, Discrete Math.\ 342 (2019) 1293--1299.
\bibitem{BS21} F.~Botler and M.~Sambinelli, Towards Gallai's path decomposition conjecture, J.~Graph
Theory 97 (2021) 161--184.
\bibitem{Ch77} B.~V.~Cherkassky, A solution of a problem on multicommodity flows in a network (in
Russian), Ekonom.\ i Mat.\ Metody 13 (1977) 143--151.
\bibitem{Edm70} J.~Edmonds, Submodular functions, matroids, and certain polyhedra, in: Combinatorial
Structures and their Applications (Calgary, 1969), Gordon and Breach, 1970, 69--87.
\bibitem{EJ73} J.~Edmonds and E.~L.~Johnson, Matching, Euler tours and the Chinese postman, Math.\
Programming 5 (1973) 88--124.
\bibitem{Fan05} G.~Fan, Path decompositions and Gallai's conjecture, J.~Combin.\ Theory Ser.~B 93 (2005)
117--125.
\bibitem{Fr92} A.~Frank, On a theorem of Mader, Discrete Math.\ 101 (1992) 49--57.
\bibitem{FrankBook} A.~Frank, Connections in Combinatorial Optimization, Oxford University Press, 2011.
\bibitem{Gale57} D.~Gale, A theorem on flows in networks, Pacific J.\ Math.\ 7 (1957) 1073--1082.
\bibitem{GJT76} M.~R.~Garey, D.~S.~Johnson and R.~E.~Tarjan, The planar Hamiltonian circuit problem is
NP-complete, SIAM J.\ Comput.\ 5 (1976) 704--714.
\bibitem{Hof60} A.~J.~Hoffman, Some recent applications of the theory of linear inequalities to extremal
combinatorial analysis, Proc.\ Sympos.\ Appl.\ Math.\ 10 (1960) 113--127.
\bibitem{Kot57} A.~Kotzig, Aus der Theorie der endlichen regul\"aren Graphen dritten und vierten Grades,
\v{C}asopis P\v{e}st.\ Mat.\ 82 (1957) 76--92.
\bibitem{L68} L.~Lov\'asz, On covering of graphs, in: Theory of Graphs (Proc.\ Colloq.\ Tihany, 1966),
Academic Press, 1968, 231--236.
\bibitem{L76} L.~Lov\'asz, On some connectivity properties of Eulerian graphs, Acta Math.\ Acad.\ Sci.\
Hungar.\ 28 (1976) 129--138.
\bibitem{Mad78} W.~Mader, A reduction method for edge-connectivity in graphs, Ann.\ Discrete Math.\ 3
(1978) 145--164.
\bibitem{OS81} H.~Okamura and P.~D.~Seymour, Multicommodity flows in planar graphs, J.~Combin.\ Theory
Ser.~B 31 (1981) 75--81.
\bibitem{Pyb96} L.~Pyber, Covering the edges of a connected graph by paths, J.~Combin.\ Theory Ser.~B 66
(1996) 152--159.
\end{thebibliography}
\end{document}
